\chapter{Fraktale Entropie – Der verhinderte Wärmetod}

\section{Einleitung}

Die Theorie der fraktalen Entropie (siehe Kapitel 13) ist eine fundamentale Erweiterung 
der WDBT+. Sie beschreibt einen geschlossenen, stationären Kreislauf der fraktalen 
Entropie zwischen Vakuum, Materie und dem $\Omega$-Feld.

Die Gravitation erzeugt durch die DSTT-Drift ($\dot{\beta} > 0$) irreversible Prozesse 
und einen intrinsischen Zeitpfeil, während die Elektrodynamik ($\dot{\beta} = 0$) 
stabile Strukturen bewahrt.

Die fraktale Raumdimension $D \approx 2,704$ verhindert ein globales Entropiemaximum 
und damit den Wärmetod.

\section{Axiome der Theorie}

\subsection{Axiom 1: Fraktale Raumstruktur}

Der fundamentale Raum ist eine fraktale Menge $M$ mit Hausdorff-Dimension $D$:

\[
D = \frac{\ln(20\phi)}{\ln(2 + \phi)} \approx 2,704, \quad \phi = \frac{1 + \sqrt{5}}{2} \tag{M.1}
\]

Die fundamentale Längenskala ist $l_0 \approx 1,8 \times 10^{-15}\,\text{m}$.

\subsection{Axiom 2: Fraktale Entropie}

Für eine Wahrscheinlichkeitsdichte $\rho: M \to \mathbb{R}_+$ mit 
$\int_M d^D r \, \rho(\vec{r}) = 1$ ist die fraktale Entropie definiert als:

\[
S_D[\rho] = -k_B \int_M d^D r \, \rho(\vec{r}) \ln \rho(\vec{r}) \cdot \left( \frac{|\vec{r}|}{l_0} \right)^{D-3} \tag{M.2}
\]

Für $D = 3$ ergibt sich die Boltzmann-Entropie. Für $D < 3$ ist $S_D$ subextensiv.

\subsection{Axiom 3: Globaler Entropieerhaltungssatz}

Die Gesamt-fraktale Entropie des Universums ist konstant:

\[
\frac{dS_D^{\text{ges}}}{dt} = 0 \tag{M.3}
\]

\subsection{Axiom 4: Lokale Entropieproduktion}

In jeder einzelnen Phase gilt der zweite Hauptsatz in erweiterter Form:

\[
\frac{dS_D^{(i)}}{dt} \geq 0 \tag{M.4}
\]

\subsection{Axiom 5: Gravitation als Quelle der Irreversibilität}

Die Weber-Gravitation mit $\beta_{\text{WG}} = 0,5$ führt zu einer monotonen Zunahme 
des tangentialen Energieanteils:

\[
\dot{\beta}_{\text{DSTT}} > 0 \quad \text{für alle } t \tag{M.5}
\]

\subsection{Axiom 6: Elektrodynamik als Stabilisator}

Die Weber-Elektrodynamik mit $\beta_{\text{WED}} = 2$ führt zu keiner Änderung der 
Energieaufteilung:

\[
\dot{\beta}_{\text{DSTT}} = 0 \quad \text{für alle } t \tag{M.6}
\]

\subsection{Axiom 7: Beschränktheit aller physikalischen Größen}

Für jede physikalische Größe $X$ (Energie, Masse, Dichte, Temperatur, Volumen, 
Komplexität) gilt:

\[
\limsup_{t \to \infty} |X(t)| < \infty \tag{M.7}
\]

Keine physikalische Größe divergiert oder wächst unbegrenzt.

\section{Die drei Phasen}

Das Universum besteht aus drei Phasen, zwischen denen die fraktale Entropie zirkuliert:

\begin{enumerate}
    \item \textbf{Vakuum} (Nullpunkt): $S_D^{\text{vac}}$ – die fraktale Struktur des Raums selbst
    \item \textbf{Materie} (kinetisch): $S_D^{\text{kin}}$ – Bewegung der Teilchen, unterteilt in:
          \begin{itemize}
              \item $S_D^{\text{kin,grav}}$ – durch Gravitation gebundene Materie
              \item $S_D^{\text{kin,em}}$ – durch Elektrodynamik gebundene Materie
          \end{itemize}
    \item \textbf{$\Omega$-Feld}: $S_D^{\Omega}$ – das Gravitationsfeld als Rotationsfeld
\end{enumerate}

Der Erhaltungssatz (M.3) lautet:

\[
S_D^{\text{vac}}(t) + S_D^{\text{kin,grav}}(t) + S_D^{\text{kin,em}}(t) + S_D^{\Omega}(t) = S_D^{\text{ges}} = \text{const} \tag{M.8}
\]

\section{Herleitung der Entropieflüsse}

\subsection{Kontinuitätsgleichung der fraktalen Entropie}

Die lokale fraktale Entropiedichte $s_D(\vec{r}, t)$ genügt einer Kontinuitätsgleichung 
mit Quellterm:

\[
\frac{\partial s_D}{\partial t} + \nabla_D \cdot \vec{j}_s = \sigma \tag{M.9}
\]

Dabei ist $\vec{j}_s$ der Entropiestrom und $\sigma$ die lokale Entropieproduktionsrate.

\subsection{Identifikation der Flussraten aus der WDBT+}

Aus den Bewegungsgleichungen der WDBT+ folgen die Flussraten durch Vergleich mit den 
Energieflüssen:

\begin{itemize}
    \item \textbf{Materieentstehung ($\sigma_{\text{cre}}$)}:  
          Das Quantenpotential $Q$ erzeugt neue Materie. Die Entropieproduktion ist 
          proportional zum Quadrat des Gradienten von $Q$:
          \[
          \sigma_{\text{cre}} = \frac{1}{T} \int \frac{|\nabla_D Q|^2}{\hbar c} \, d^D r \tag{M.10}
          \]
    
    \item \textbf{DSTT-Drift ($\sigma_{\text{drift}}$)}:  
          Die Drift von $\beta_{\text{DSTT}}$ führt zu einem Entropiefluss in das 
          $\Omega$-Feld:
          \[
          \sigma_{\text{drift}} = \frac{1}{T} \int \frac{\dot{\beta}_{\text{DSTT}}}{\beta_{\text{DSTT}}} \, d^D r \tag{M.11}
          \]
    
    \item \textbf{Rückkopplung ($\sigma_{\text{back}}$)}:  
          Die Energie des $\Omega$-Feldes wird an das Vakuum zurückgegeben:
          \[
          \sigma_{\text{back}} = \frac{1}{T} \int \frac{|\vec{\Omega}|^2}{8\pi G} \, d^D r \tag{M.12}
          \]
\end{itemize}

\subsection{Positivität der Flussraten}

Alle drei Flussraten sind positiv definit:
\begin{itemize}
    \item $|\nabla_D Q|^2 > 0$ für nicht-triviale Q-Felder
    \item $\dot{\beta}_{\text{DSTT}} > 0$ aus der DSTT-Drift
    \item $|\vec{\Omega}|^2 > 0$ für nicht-verschwindende $\Omega$-Felder
\end{itemize}

\subsection{Entropiebilanz der einzelnen Phasen}

Die zeitlichen Änderungen der Entropien jeder Phase ergeben sich aus den Flussraten:

\begin{align}
\frac{dS_D^{\text{vac}}}{dt} &= -\sigma_{\text{cre}} + \sigma_{\text{back}} \tag{M.13a} \\
\frac{dS_D^{\text{kin,grav}}}{dt} &= +\sigma_{\text{cre}} - \sigma_{\text{drift}} \tag{M.13b} \\
\frac{dS_D^{\text{kin,em}}}{dt} &= 0 \tag{M.13c} \\
\frac{dS_D^{\Omega}}{dt} &= +\sigma_{\text{drift}} - \sigma_{\text{back}} \tag{M.13d}
\end{align}

\section{Der geschlossene Kreislauf}

\subsection{Darstellung des Kreislaufs}

Der geschlossene Entropiekreislauf lässt sich wie folgt darstellen:

\[
\text{Vakuum} \xrightarrow{\sigma_{\text{cre}}} \text{Materie (grav.)} \xrightarrow{\sigma_{\text{drift}}} \Omega\text{-Feld} \xrightarrow{\sigma_{\text{back}}} \text{Vakuum}
\]

Die elektromagnetische Materie ist vom Kreislauf entkoppelt (reversibel).

\subsection{Stationärer Zustand}

Im stationären Zustand gilt:

\[
\sigma_{\text{cre}} = \sigma_{\text{drift}} = \sigma_{\text{back}} \equiv \sigma \tag{M.14}
\]

Damit sind die Entropien jeder Phase konstant:

\[
\frac{dS_D^{\text{vac}}}{dt} = \frac{dS_D^{\text{kin,grav}}}{dt} = \frac{dS_D^{\text{kin,em}}}{dt} = \frac{dS_D^{\Omega}}{dt} = 0 \tag{M.15}
\]

Es zirkuliert nur der Fluss, nicht die Menge.

\section{Warum kein Wärmetod eintritt}

\subsection{Bedingungen für einen Wärmetod}

Ein Wärmetod tritt ein, wenn:
\begin{enumerate}
    \item Die Entropie ein globales Maximum erreicht
    \item Alle Energiegradienten verschwinden
    \item Keine irreversiblen Prozesse mehr stattfinden
    \item Das System in einen Gleichgewichtszustand übergeht
\end{enumerate}

\subsection{Verhinderung des Wärmetods durch fraktale Geometrie}

Die fraktale Entropie $S_D$ hat bei fester Gesamtenergie kein globales Maximum. 
Die Zustandsdichte im fraktalen Raum skaliert mit $E^{D/3}$. Für $D < 3$ ist die 
Funktion

\[
S_D(E) = \frac{D}{3}k_B \ln E + \text{const} \tag{M.16}
\]

monoton wachsend und konkav, besitzt aber kein endliches Maximum.

\subsection{Beweis: Kein endliches Entropiemaximum}

Angenommen, es gäbe ein endliches Maximum bei $E = E_0$. Dann müsste für $E > E_0$ 
gelten: $S_D(E) < S_D(E_0)$. Aus (M.16) folgt jedoch für $E > E_0$:

\[
S_D(E) - S_D(E_0) = \frac{D}{3}k_B \ln \left( \frac{E}{E_0} \right) > 0 \quad \text{für } E > E_0 \tag{M.17}
\]

Widerspruch. Also existiert kein endliches Maximum.

\subsection{Permanente Gradienten}

Die DSTT-Drift ($\dot{\beta} > 0$) erzeugt permanent radiale Energiegradienten:

\[
\nabla E_{\text{kin,grav}} \neq 0 \quad \text{für alle } t \tag{M.18}
\]

Die Raten $\sigma_{\text{cre}}, \sigma_{\text{drift}}, \sigma_{\text{back}}$ sind immer positiv.

\subsection{Hauptsatz der fraktalen Thermodynamik}

Unter den Axiomen 1–7 gilt:

\begin{enumerate}
    \item \textbf{Lokale Entropieproduktion:} $\displaystyle \frac{dS_D^{(i)}}{dt} \geq 0$ für jede Phase
    \item \textbf{Globale Entropieerhaltung:} $\displaystyle \frac{dS_D^{\text{ges}}}{dt} = 0$
    \item \textbf{Permanente Dynamik:} $\sigma_{\text{cre}}, \sigma_{\text{drift}}, \sigma_{\text{back}} > 0$ für alle $t$
    \item \textbf{Kein Entropiemaximum:} Es existiert kein endlicher Zustand maximaler Gesamtentropie
    \item \textbf{Kein Wärmetod:} Das System erreicht nie einen Gleichgewichtszustand mit verschwindenden Gradienten und Prozessen
\end{enumerate}

\section{Kosmologische Konsequenzen}

Die Theorie der fraktalen Entropie führt zu einem stationären Universum:

\begin{itemize}
    \item Der Raum expandiert nicht, aber Strukturen und Bahnen dehnen sich durch die 
          DSTT-Drift aus.
    \item Die mittlere Dichte bleibt durch kontinuierliche Materieentstehung (Anhang J) konstant.
    \item Die Gravitation produziert permanent Entropie – dies ist der intrinsische Zeitpfeil.
    \item Die Elektrodynamik ermöglicht reversible Prozesse und bewahrt stabile Strukturen.
    \item Der Wärmetod tritt nicht ein – das Universum bleibt ewig dynamisch.
\end{itemize}

\subsection{Temperatur im stationären Zustand}

Die mittlere Temperatur des Universums folgt aus der Gleichgewichtsbedingung zwischen 
Entropieproduktion und -abfluss:

\[
T = \left( \frac{\partial S_D}{\partial E} \right)^{-1} = \frac{3}{D} \frac{E}{k_B} \tag{M.19}
\]

Für $D < 3$ ist die effektive Temperatur höher als im $D = 3$-Fall.

\subsection{Zeitpfeil und kosmologische Entropie}

Die totale Entropieproduktion des Universums ist:

\[
\dot{S}_D^{\text{ges}} = \dot{S}_D^{\text{grav}} + \dot{S}_D^{\text{em}} + \dot{S}_D^{\text{vac}} = \dot{S}_D^{\text{grav}} > 0 \tag{M.20}
\]

Die Gravitation bricht die Zeitumkehrsymmetrie. Dies ist der fundamentale Ursprung 
des kosmologischen Zeitpfeils – nicht die Anfangsbedingungen des Urknalls, sondern 
die intrinsische Irreversibilität der Gravitation selbst.

\section{Zusammenfassung}

\begin{itemize}
    \item Der Raum ist fraktal mit $D \approx 2,704$ und besitzt eine fundamentale Länge $l_0$.
    \item Die fraktale Entropie $S_D$ ist subextensiv und zirkuliert in einem geschlossenen 
          Kreislauf zwischen Vakuum, Materie und $\Omega$-Feld.
    \item Die Entropieflüsse werden aus den Bewegungsgleichungen der WDBT+ hergeleitet:
          $\sigma_{\text{cre}} \propto |\nabla_D Q|^2$, 
          $\sigma_{\text{drift}} \propto \dot{\beta}/\beta$, 
          $\sigma_{\text{back}} \propto |\vec{\Omega}|^2$.
    \item Die Gravitation ($\beta_{\text{WG}} = 0,5$) erzeugt durch $\dot{\beta} > 0$ 
          irreversible Prozesse und den intrinsischen Zeitpfeil.
    \item Die Elektrodynamik ($\beta_{\text{WED}} = 2$) mit $\dot{\beta} = 0$ ermöglicht 
          reversible Prozesse und bewahrt stabile Strukturen.
    \item Der Wärmetod wird durch die fraktale Geometrie (kein globales Entropiemaximum) 
          und die permanenten irreversiblen Flüsse verhindert.
    \item Das Universum ist stationär: Energie, Masse und Dichte bleiben konstant; 
          die Dynamik zirkuliert ewig.
\end{itemize}